\documentclass[11pt]{letter}
\usepackage[letterpaper,margin=0.9in]{geometry}
\usepackage[hidelinks]{hyperref}
\usepackage{parskip}
\pagenumbering{gobble}

\begin{document}

\begin{center}
{\Large \textbf{Aryan Miriyala}}\\
\href{mailto:aryanmiriyala@gmail.com}{aryanmiriyala@gmail.com} $|$ +1 419-315-0444 $|$
\href{https://www.linkedin.com/in/aryan-miriyala-7788a21b6/}{linkedin.com/in/aryan-miriyala} $|$
\href{https://github.com/aryanmiriyala}{github.com/aryanmiriyala}
\end{center}

\vspace{0.25in}

Dear GlobalFoundries Hiring Team,

GlobalFoundries' AI/ML Systems Engineer role is interesting to me because it sits exactly where I have been trying to grow: the boundary between machine learning workloads, systems behavior, and quantitative performance reasoning. I am drawn to the idea of helping a semiconductor team understand how AI workloads behave on real or projected hardware, especially at a company whose work connects essential chips, power-efficient technologies, and AI infrastructure from the cloud to the physical world. My background is not in silicon design, so I would not pretend otherwise, but I do bring a strong computer science foundation, a 4.00 GPA through my B.S. and M.S. work, and practical systems and data projects where the job was to measure behavior carefully, build reproducible analysis, and explain the results clearly.

The closest example is my self-adaptive OpenMP scheduling project. I built a C++/OpenMP runtime experiment that collected per-epoch telemetry and used a UCB-based controller to tune scheduling policy, chunk size, and thread count, then benchmarked it against serial and fixed-policy baselines across Mandelbrot, heat diffusion, and reduction workloads. I generated plots for runtime, speedup, efficiency, Karp-Flatt behavior, and scheduler decisions, which forced me to think carefully about workload characteristics instead of treating performance as a single number. I have also built Python/PySpark production workflows over 20+ TB of insurance data, Python/Pandas data-cleaning pipelines, and ML/RL experiments with PyTorch, Stable-Baselines3, and TensorBoard. I would bring GF a careful early-career engineer's mindset: write clean analysis code, validate assumptions, communicate tradeoffs clearly, and keep learning from architecture, compiler, runtime, and hardware specialists around me.

\vspace{0.15in}

Sincerely,\\
Aryan Miriyala

\end{document}
